\documentclass[11pt]{article}

% Shared notes settings for Elements of AIML lectures (UPES).
% Keep this file free of \documentclass to allow reuse via \input.

\usepackage[utf8]{inputenc}
\usepackage[T1]{fontenc}
\usepackage[a4paper,margin=1in]{geometry}
\usepackage{hyperref}
\usepackage{xurl}
\usepackage{booktabs}
\usepackage{amsmath}
\usepackage{amssymb}
\usepackage{listings}
\usepackage{xcolor}
\usepackage{enumitem}
\usepackage{parskip}

% Code listings (general-purpose).
\lstset{
  basicstyle=\ttfamily\small,
  keywordstyle=\color{blue},
  commentstyle=\color{gray},
  stringstyle=\color{teal},
  showstringspaces=false,
  columns=fullflexible,
  frame=single,
  framerule=0.2pt,
  breaklines=true
}

\setlist[itemize]{topsep=0.3em,itemsep=0.2em}
\setlist[enumerate]{topsep=0.3em,itemsep=0.2em}

% Simple per-section source line.
\newcommand{\notesources}[1]{\par\noindent{\small\color{gray}\textit{Sources:} \sloppy #1\par}}

% Unit-wise source strings for this subject (used by slides/notes).
% Path is relative to the lecture `latex/` folder (the compilation CWD).
\IfFileExists{../../../_shared/aiml-sources.tex}{%
  % Source strings for Elements of AIML (used by slides + notes).

\newcommand{\AIMLRepoURL}{https://github.com/tali7c/Elements-of-AIML}

\newcommand{\AIMLLabSources}{%
Provost and Fawcett, \textit{Data Science for Business}.;%
McKinney, \textit{Python for Data Analysis}.;%
WEKA Documentation (Waikato Environment for Knowledge Analysis), accessed 2026-02-28.;%
Matplotlib and Seaborn documentation, accessed 2026-02-28.%
}

\newcommand{\AIMLUnitISources}{%
Rich and Knight, \textit{Artificial Intelligence}, McGraw Hill, 2017.;%
Russell and Norvig, \textit{Artificial Intelligence: A Modern Approach} (4e), Ch. 1--2 (Introduction, Intelligent Agents).;%
Poole and Mackworth, \textit{Artificial Intelligence: Foundations of Computational Agents} (2e), selected sections.%
}

\newcommand{\AIMLUnitIISources}{%
Russell and Norvig, \textit{Artificial Intelligence: A Modern Approach} (4e), Ch. 7--9 (Logical Agents, First-Order Logic, Inference).;%
Huth and Ryan, \textit{Logic in Computer Science} (2e), selected sections on propositional and first-order logic.;%
Genesereth and Nilsson, \textit{Logical Foundations of Artificial Intelligence}, selected chapters.%
}

\newcommand{\AIMLUnitIIISources}{%
Mueller and Massaron, \textit{Machine Learning for Dummies}, For Dummies, 2016.;%
Mitchell, \textit{Machine Learning}, selected chapters on learning basics and evaluation.;%
Scikit-learn documentation, accessed 2026-02-28.%
}

\newcommand{\AIMLUnitIVSources}{%
Mueller and Massaron, \textit{Machine Learning for Dummies}, For Dummies, 2016.;%
James, Witten, Hastie, and Tibshirani, \textit{An Introduction to Statistical Learning} (2e), selected chapters.;%
Scikit-learn documentation, accessed 2026-02-28.%
}

\newcommand{\AIMLUnitVSources}{%
NIST, \textit{AI Risk Management Framework (AI RMF 1.0)}, 2023.;%
Barocas, Hardt, and Narayanan, \textit{Fairness and Machine Learning} (online book), selected sections.;%
Knaflic, \textit{Storytelling with Data}, selected chapters on communicating results.%
}
%
}{%
  % Faculty copies live under `private/` so their compile CWD is different.
  \IfFileExists{../../../../public/_shared/aiml-sources.tex}{%
    % Source strings for Elements of AIML (used by slides + notes).

\newcommand{\AIMLRepoURL}{https://github.com/tali7c/Elements-of-AIML}

\newcommand{\AIMLLabSources}{%
Provost and Fawcett, \textit{Data Science for Business}.;%
McKinney, \textit{Python for Data Analysis}.;%
WEKA Documentation (Waikato Environment for Knowledge Analysis), accessed 2026-02-28.;%
Matplotlib and Seaborn documentation, accessed 2026-02-28.%
}

\newcommand{\AIMLUnitISources}{%
Rich and Knight, \textit{Artificial Intelligence}, McGraw Hill, 2017.;%
Russell and Norvig, \textit{Artificial Intelligence: A Modern Approach} (4e), Ch. 1--2 (Introduction, Intelligent Agents).;%
Poole and Mackworth, \textit{Artificial Intelligence: Foundations of Computational Agents} (2e), selected sections.%
}

\newcommand{\AIMLUnitIISources}{%
Russell and Norvig, \textit{Artificial Intelligence: A Modern Approach} (4e), Ch. 7--9 (Logical Agents, First-Order Logic, Inference).;%
Huth and Ryan, \textit{Logic in Computer Science} (2e), selected sections on propositional and first-order logic.;%
Genesereth and Nilsson, \textit{Logical Foundations of Artificial Intelligence}, selected chapters.%
}

\newcommand{\AIMLUnitIIISources}{%
Mueller and Massaron, \textit{Machine Learning for Dummies}, For Dummies, 2016.;%
Mitchell, \textit{Machine Learning}, selected chapters on learning basics and evaluation.;%
Scikit-learn documentation, accessed 2026-02-28.%
}

\newcommand{\AIMLUnitIVSources}{%
Mueller and Massaron, \textit{Machine Learning for Dummies}, For Dummies, 2016.;%
James, Witten, Hastie, and Tibshirani, \textit{An Introduction to Statistical Learning} (2e), selected chapters.;%
Scikit-learn documentation, accessed 2026-02-28.%
}

\newcommand{\AIMLUnitVSources}{%
NIST, \textit{AI Risk Management Framework (AI RMF 1.0)}, 2023.;%
Barocas, Hardt, and Narayanan, \textit{Fairness and Machine Learning} (online book), selected sections.;%
Knaflic, \textit{Storytelling with Data}, selected chapters on communicating results.%
}
%
  }{%
    % Source strings for Elements of AIML (used by slides + notes).

\newcommand{\AIMLRepoURL}{https://github.com/tali7c/Elements-of-AIML}

\newcommand{\AIMLLabSources}{%
Provost and Fawcett, \textit{Data Science for Business}.;%
McKinney, \textit{Python for Data Analysis}.;%
WEKA Documentation (Waikato Environment for Knowledge Analysis), accessed 2026-02-28.;%
Matplotlib and Seaborn documentation, accessed 2026-02-28.%
}

\newcommand{\AIMLUnitISources}{%
Rich and Knight, \textit{Artificial Intelligence}, McGraw Hill, 2017.;%
Russell and Norvig, \textit{Artificial Intelligence: A Modern Approach} (4e), Ch. 1--2 (Introduction, Intelligent Agents).;%
Poole and Mackworth, \textit{Artificial Intelligence: Foundations of Computational Agents} (2e), selected sections.%
}

\newcommand{\AIMLUnitIISources}{%
Russell and Norvig, \textit{Artificial Intelligence: A Modern Approach} (4e), Ch. 7--9 (Logical Agents, First-Order Logic, Inference).;%
Huth and Ryan, \textit{Logic in Computer Science} (2e), selected sections on propositional and first-order logic.;%
Genesereth and Nilsson, \textit{Logical Foundations of Artificial Intelligence}, selected chapters.%
}

\newcommand{\AIMLUnitIIISources}{%
Mueller and Massaron, \textit{Machine Learning for Dummies}, For Dummies, 2016.;%
Mitchell, \textit{Machine Learning}, selected chapters on learning basics and evaluation.;%
Scikit-learn documentation, accessed 2026-02-28.%
}

\newcommand{\AIMLUnitIVSources}{%
Mueller and Massaron, \textit{Machine Learning for Dummies}, For Dummies, 2016.;%
James, Witten, Hastie, and Tibshirani, \textit{An Introduction to Statistical Learning} (2e), selected chapters.;%
Scikit-learn documentation, accessed 2026-02-28.%
}

\newcommand{\AIMLUnitVSources}{%
NIST, \textit{AI Risk Management Framework (AI RMF 1.0)}, 2023.;%
Barocas, Hardt, and Narayanan, \textit{Fairness and Machine Learning} (online book), selected sections.;%
Knaflic, \textit{Storytelling with Data}, selected chapters on communicating results.%
}
%
  }%
}


\title{Elements of AIML\\Miscellaneous -- Lecture 01 Notes\\PCA, LDA, and ICA in One View}
\author{Tofik Ali}
\date{\today}

\begin{document}
\maketitle
\tableofcontents

\section{Lecture Overview}
This supplementary lecture puts \textbf{PCA}, \textbf{LDA}, and \textbf{ICA} side by side.
Students often remember their names but confuse:
\begin{itemize}[leftmargin=*]
  \item what each method optimizes,
  \item whether labels are needed,
  \item what kind of output each method produces,
  \item when one method should be preferred over another.
\end{itemize}
The aim here is not to replace the Unit III lecture, but to make the comparison easier to revise later.
\notesources{\AIMLMiscSources}

\section{How to use these notes}
Recommended order:
\begin{enumerate}[leftmargin=*]
  \item First read \textbf{Unit 03 Lecture 03} for PCA basics and the leakage-safe preprocessing rule.
  \item Read this lecture to compare PCA, LDA, and ICA more systematically.
  \item Read \textbf{Miscellaneous Lecture 02} if you want the full manual eigenvalue/eigenvector derivation behind PCA.
\end{enumerate}

\section{What should already be familiar?}
\begin{itemize}[leftmargin=*]
  \item From \textbf{Unit 03 Lecture 03}: covariance matrix, eigenvalues, eigenvectors, and explained variance.
  \item From \textbf{Unit 04}: class labels, classes, and why classification quality depends on separability.
\end{itemize}
If these ideas are weak, revisit those lectures first; otherwise the comparison here will feel more abstract than useful.

\section{A common confusion: same category, different goals}
PCA, LDA, and ICA are all discussed under \textbf{dimensionality reduction} or \textbf{feature transformation}, but they do not solve the same problem.

\subsection{A useful first question}
Before choosing one of the three methods, ask:
\begin{enumerate}[leftmargin=*]
  \item Do I have labels?
  \item Do I want compression, class separation, or latent-source recovery?
  \item Do I care more about variance, discrimination, or independence?
\end{enumerate}

\section{PCA}
\subsection{Core idea}
\textbf{Principal Component Analysis (PCA)} finds orthogonal directions that capture the largest variance in the data.
It is unsupervised: labels are not used.

\subsection{What PCA optimizes}
If $C$ is the covariance matrix, PCA chooses a unit vector $v$ that maximizes projected variance:
\[
  \max_{\|v\|=1} v^\top C v
\]
The first solution is the eigenvector corresponding to the largest eigenvalue.
The second solution is orthogonal to the first, and so on.

\subsection{What PCA gives back}
\begin{itemize}[leftmargin=*]
  \item principal directions (eigenvectors),
  \item variance captured by each direction (eigenvalues),
  \item transformed coordinates in the reduced space.
\end{itemize}

\subsection{When PCA is useful}
\begin{itemize}[leftmargin=*]
  \item data compression,
  \item noise reduction,
  \item 2D or 3D visualization,
  \item reducing redundancy in correlated features.
\end{itemize}

\subsection{Where PCA can disappoint}
\begin{itemize}[leftmargin=*]
  \item High variance is not always the same as high usefulness.
  \item PCA may keep directions that are statistically strong but not class-discriminative.
  \item Components are often harder to interpret than original features.
\end{itemize}

\section{LDA}
\subsection{Core idea}
\textbf{Linear Discriminant Analysis (LDA)} tries to find projections that separate classes well.
Unlike PCA, it is supervised because it uses class labels.

\subsection{The two scatter matrices}
LDA balances:
\begin{itemize}[leftmargin=*]
  \item \textbf{within-class scatter} $S_W$: how spread each class is inside itself,
  \item \textbf{between-class scatter} $S_B$: how far class means are from each other.
\end{itemize}
The usual criterion is:
\[
  J(W) = \frac{|W^\top S_B W|}{|W^\top S_W W|}
\]
So the method wants:
\begin{itemize}[leftmargin=*]
  \item large separation between class centers,
  \item small spread within each class.
\end{itemize}

\subsection{What LDA gives back}
LDA returns discriminant directions.
These directions are chosen for class separation, not for overall data variance.

\subsection{Why LDA has a dimension limit}
If there are $C$ classes, LDA can return at most $C-1$ useful directions.
This happens because between-class information depends on class means, and there are only $C-1$ independent mean-separation directions.

\subsection{When LDA is useful}
\begin{itemize}[leftmargin=*]
  \item when class labels are available,
  \item when the purpose is classification support,
  \item when you want a lower-dimensional representation that still preserves class discrimination.
\end{itemize}

\subsection{Where LDA can disappoint}
\begin{itemize}[leftmargin=*]
  \item If classes overlap strongly, LDA cannot force clean separation.
  \item If labels are noisy, the projection becomes less meaningful.
  \item Without labels, LDA cannot be applied at all.
\end{itemize}

\section{ICA}
\subsection{Core idea}
\textbf{Independent Component Analysis (ICA)} assumes observed features are mixtures of hidden latent sources.
The aim is to recover components that are as statistically independent as possible.

\subsection{Mixing model}
ICA is usually explained by:
\[
  X = A S
\]
where:
\begin{itemize}[leftmargin=*]
  \item $S$ contains hidden source signals,
  \item $A$ is an unknown mixing matrix,
  \item $X$ is what we observe.
\end{itemize}
The job of ICA is to estimate an unmixing matrix so that the recovered components behave independently.

\subsection{Why ICA is different from PCA}
PCA only removes correlation and ranks directions by variance.
ICA goes further: it tries to recover statistically independent components, which is a stronger condition than uncorrelatedness.

\subsection{Typical ICA pipeline}
\begin{enumerate}[leftmargin=*]
  \item center the data,
  \item often whiten the data,
  \item optimize an independence criterion, often related to non-Gaussianity.
\end{enumerate}

\subsection{When ICA is useful}
\begin{itemize}[leftmargin=*]
  \item signal separation,
  \item EEG or audio artifact removal,
  \item recovering hidden source-like factors from mixed observations.
\end{itemize}

\subsection{Where ICA can disappoint}
\begin{itemize}[leftmargin=*]
  \item If the latent sources are not independent, the model assumption breaks.
  \item If the data is close to Gaussian, ICA becomes much less informative.
  \item ICA components are often sensitive to scaling and preprocessing choices.
\end{itemize}

\section{Detailed comparison}
\subsection{Comparison by objective}
\begin{center}
\begin{tabular}{p{2.2cm}p{5.2cm}p{5.2cm}}
\toprule
\textbf{Method} & \textbf{Main objective} & \textbf{Interpretation} \\
\midrule
PCA & maximize variance & keep the strongest spread directions \\
LDA & maximize between-class separation relative to within-class spread & help classes separate in a lower-dimensional space \\
ICA & maximize statistical independence & recover hidden source-like signals \\
\bottomrule
\end{tabular}
\end{center}

\subsection{Comparison by supervision}
\begin{center}
\begin{tabular}{p{2.2cm}p{2.5cm}p{7.9cm}}
\toprule
\textbf{Method} & \textbf{Uses labels?} & \textbf{Why this matters} \\
\midrule
PCA & No & cannot use label information, so class separation may be accidental rather than intentional \\
LDA & Yes & uses labels directly, so it can align the representation with classification needs \\
ICA & No & focuses on independence structure, not on labels or class boundaries \\
\bottomrule
\end{tabular}
\end{center}

\subsection{Comparison by assumptions}
\begin{itemize}[leftmargin=*]
  \item \textbf{PCA}: important information is reflected in large-variance directions.
  \item \textbf{LDA}: classes are meaningful and separability matters.
  \item \textbf{ICA}: latent sources are independent and often non-Gaussian.
\end{itemize}

\subsection{Comparison by geometry/statistics}
\begin{itemize}[leftmargin=*]
  \item \textbf{PCA} produces orthogonal components.
  \item \textbf{LDA} produces class-discriminative directions, not necessarily variance-dominant ones.
  \item \textbf{ICA} seeks independence, which is stronger than orthogonality and stronger than simple uncorrelatedness.
\end{itemize}

\subsection{Comparison by output dimension}
\begin{itemize}[leftmargin=*]
  \item \textbf{PCA}: any $k \le d$ may be chosen.
  \item \textbf{LDA}: at most $C-1$ directions for $C$ classes.
  \item \textbf{ICA}: usually chosen according to the assumed number of sources or observed variables.
\end{itemize}

\subsection{Comparison by classroom use-case}
\begin{itemize}[leftmargin=*]
  \item If you want to compress features before plotting or faster modeling, start with \textbf{PCA}.
  \item If your real goal is classification and labels are available, \textbf{LDA} may be more aligned with the task.
  \item If you are separating mixed signals or hidden factors, \textbf{ICA} is the more natural tool.
\end{itemize}

\section{One-sentence memory anchors}
\begin{itemize}[leftmargin=*]
  \item \textbf{PCA:} ``keep the directions with the most variance.''
  \item \textbf{LDA:} ``keep the directions that separate classes.''
  \item \textbf{ICA:} ``recover independent hidden sources from mixtures.''
\end{itemize}

\section{What students often say incorrectly}
\subsection*{Incorrect statement 1: ``LDA is supervised PCA.''}
\textbf{Correction:} the two methods are related only at a high level as linear transformations. Their objectives are different: PCA maximizes variance, while LDA maximizes discriminative separation.

\subsection*{Incorrect statement 2: ``ICA gives more important principal components.''}
\textbf{Correction:} ICA does not rank components by variance at all. It looks for independence, not for variance dominance.

\subsection*{Incorrect statement 3: ``Orthogonal means independent.''}
\textbf{Correction:} orthogonal vectors are geometrically perpendicular; independence is a stronger probabilistic statement. Two components can be uncorrelated but still not independent.

\section{Decision checklist}
Use this quick decision process:
\begin{enumerate}[leftmargin=*]
  \item If labels are absent, remove LDA from the candidate list.
  \item If the goal is compression or visualization, prefer PCA.
  \item If the goal is better class separation, investigate LDA.
  \item If the goal is source separation or mixed-signal recovery, investigate ICA.
  \item Check whether the method assumptions actually make sense for the data.
\end{enumerate}

\section{Practice Questions (with answers)}
\subsection*{Q1. Why can PCA fail to help classification even when it keeps most variance?}
\textbf{Answer:} because the directions that contain the most variance may not be the directions that best separate classes. Variance and discriminative power are different objectives.

\subsection*{Q2. Why can LDA not be used without labels?}
\textbf{Answer:} because LDA needs class means and class-wise scatter. Without labels, those class-based quantities are undefined.

\subsection*{Q3. Why is ICA more demanding than PCA?}
\textbf{Answer:} PCA only uses second-order structure such as covariance, while ICA tries to recover statistically independent latent sources, which requires stronger assumptions and harder optimization.

\subsection*{Q4. Which method would you try first for a 2D visualization of high-dimensional embeddings?}
\textbf{Answer:} PCA, because its goal is compression/variance preservation and it is easy to apply without labels.

\section{Summary}
PCA, LDA, and ICA belong to the same broad family of linear transformations, but they answer different questions:
\begin{itemize}[leftmargin=*]
  \item \textbf{PCA:} what directions preserve the most variance?
  \item \textbf{LDA:} what directions separate known classes best?
  \item \textbf{ICA:} what hidden components are most independent?
\end{itemize}
The method should therefore be chosen by \textbf{goal and assumptions}, not by familiarity of name.
\notesources{\AIMLMiscSources}

\end{document}
